\section{Introduction}

Scientific computing generates data at unprecedented scale: a single exascale job on Frontier can produce 80~PB of simulation output. Yet an estimated 60 to 80\% of scientific data is never reused after creation, not because it is inaccessible, but because it is unsearchable. The FAIR data principles, now mandated by NSF, NIH, and DOE, require that data be Findable, Accessible, Interoperable, and Reusable. Automated FAIR assessments show that no major repository achieves even 50\% compliance on reusability, with measurement heterogeneity a systematic barrier. Nature's survey found that over 70\% of researchers have failed to reproduce others' experiments~\cite{baker2016reproducibility}. The problem is not storage or access. It is search.

This gap is intensifying as AI agents enter scientific workflows. Coscientist~\cite{coscientist2023} autonomously executes wet-lab synthesis. AI Scientist v2~\cite{aiscientistv2_2025} generates research end-to-end. Autonomous agents at ORNL orchestrate experiments across facilities~\cite{ornl_agents_sc25}. These agents can reason about science, plan multi-step experiments, and call external tools. What they cannot do is search scientific data effectively, because the retrieval infrastructure they depend on treats all content as plain text. ScienceAgentBench~\cite{scienceagentbench2025} reports 32.4\% end-to-end task success, with failures traced to data handling. A comprehensive survey~\cite{agentic_ai_survey2025} identifies scientific data management as explicitly underexplored in the agent literature.

Three compounding failures explain why retrieval systems cannot support scientific data discovery, whether by human scientists or autonomous agents.

\textbf{Dimensional blindness.} Scientific measurements are routinely expressed in different unit prefixes. Pressure alone appears as Pa, hPa, kPa, MPa, bar, atm, and psi across disciplines. A query for 200~kPa should match ``200000~Pa'' in an HDF5 file and ``0.2~MPa'' in an S3 dataset, but no retrieval operator performs the required arithmetic. String normalization~\cite{numbersmatter2024} maps unit names to canonical forms but cannot bridge prefix gaps. Learned embeddings~\cite{cone2026} achieve 0.54 binary retrieval accuracy on numerical content, effectively random~\cite{numeracygap2026}. Unit confusion has caused real-world disasters: NASA's \$327M Mars Climate Orbiter loss (pound-seconds vs.\ newton-seconds), the Gimli Glider incident (kg vs.\ lb fuel loading), and an estimated 7,000 to 9,000 medication deaths per year from unit errors in US healthcare.

\textbf{Formula opacity.} ``$F = ma$,'' ``$F = m \cdot a$,'' and ``$ma = F$'' express the same physical law but produce distinct token sequences. Math-aware systems~\cite{approach0, ssemb2025} match equation structure but cannot connect a matched formula to the physical measurements it governs. No retrieval system combines formula normalization with quantity-aware dimensional search.

\textbf{Storage fragmentation.} HPC data spans parallel filesystems, S3 object stores, vector databases, and binary formats such as HDF5 and NetCDF. No system provides unified search with science-aware operators across these backends.

We present CLIO Search, an agentic retrieval system that addresses all three failures. The core design principle is to separate search reasoning from search execution. Unlike conventional retrieval systems that apply a fixed pipeline to every query, CLIO reasons about how to search before searching. It profiles each corpus to discover what metadata exists, selects retrieval branches based on corpus characteristics (skipping branches that cannot contribute), executes science-aware operators as deterministic database predicates alongside BM25 and dense vector search, and federates across heterogeneous storage backends. An iterative refinement loop expands, narrows, or pivots the query based on intermediate results. Domain knowledge enters the pipeline as executable, verifiable computation: the fact that 5~km = 5000~m is an algebraic invariant encoded in the retrieval operator, not a statistical similarity learned from training data.

We make five contributions:

\begin{enumerate}
\item \textbf{Dimensional-conversion retrieval.} A composable unit registry based on dimensional analysis canonicalizes measurements to SI base units at index time and evaluates cross-unit queries as database range predicates. The registry covers 46 units across 12 domains with correctness guaranteed by construction.

\item \textbf{Formula normalization.} Deterministic normalization matches structurally equivalent equations regardless of surface notation, integrated as a parallel retrieval branch.

\item \textbf{Corpus profiling and adaptive search.} A profiling layer inspects each namespace before querying, enabling metadata-adaptive strategy selection and intelligent branch selection that skips retrieval paths unable to contribute results.

\item \textbf{Federated multi-backend search.} A connector architecture supports unified retrieval across filesystem, S3, HDF5, NetCDF, IOWarp, NDP, and other backends with per-connector capability negotiation.

\item \textbf{Evaluation at multiple scales.} Controlled benchmarks (210 documents, 80 queries), real observatory data (288 NOAA documents), a live external catalog (341 NDP datasets), and distributed HPC infrastructure (70K arXiv documents across 4 GH200 nodes on NCSA DeltaAI).
\end{enumerate}

The remainder of this paper is organized as follows. Section~2 surveys hybrid retrieval, the numeracy crisis, quantity-aware retrieval, and scientific data discovery. Section~3 presents the system design including the agentic orchestration layer. Section~4 describes the implementation. Section~5 evaluates across dimensional conversion, agentic efficiency, federated search, and distributed scaling. Section~6 concludes.

%% ===================================================================
%% Section 2: Background and Related Work
